\documentclass[12pt]{article}
\usepackage[ukrainian]{babel}
\usepackage{fontspec}
% --- НАЛАШТУВАННЯ ШРИФТІВ ---
\setmainfont{Schoolbook-Regular.otf}[
    Path = ../,
    BoldFont = Schoolbook-Bold.otf,
    ItalicFont = Schoolbook-Italic.otf,
    BoldItalicFont = Schoolbook-BoldItalic.otf
]
\usepackage[italic]{mathastext}
\usepackage[a4paper,margin=1.8cm,bottom=2cm,top=2cm]{geometry}
\usepackage{amsmath,amssymb}
\usepackage{tikz}
\usepackage{pgfplots}
\pgfplotsset{compat=1.16}
\usetikzlibrary{calc,patterns,angles,quotes,intersections}
\usepackage{xcolor,array,fancyhdr,enumitem,multicol}
\usepackage{colortbl}
\usepackage{tcolorbox}
\tcbuselibrary{skins,breakable}
\usepackage[hidelinks]{hyperref}
\usepackage[firstpage=false, text=@pvtr2525, scale=0.6, color=gray!12]{draftwatermark}

% Заборона розриву інлайн-формул
\binoppenalty=10000
\relpenalty=10000

\definecolor{mainGreen}{RGB}{34, 120, 64}
\definecolor{yearOrange}{RGB}{220, 100, 30}
\definecolor{titleBg}{RGB}{225, 240, 220}
\definecolor{instrBg}{RGB}{255, 240, 220}
\definecolor{instrBorder}{RGB}{220, 150, 80}
\definecolor{gold}{RGB}{240, 190, 40}
\definecolor{silver}{RGB}{170, 170, 175}
\definecolor{bronze}{RGB}{190, 120, 70}

\pagestyle{fancy}
\fancyhf{}
\renewcommand{\headrulewidth}{0.4pt}
\renewcommand{\footrulewidth}{0.4pt}
\fancyhead[C]{\small\color{gray!80} Збірник НМТ \textendash{} Варіант за 15 червня}
\fancyhead[R]{\small\color{mainGreen}\textbf{\href{https://t.me/pvtr2525}{@pvtr2525}}}
\fancyfoot[L]{\small\color{gray!80}\href{https://t.me/pvtr2525}{ТГ: @pvtr2525}}
\fancyfoot[C]{\small\color{gray!80}\thepage}
\fancyfoot[R]{\small\color{gray!80}НМТ 2026}

\newcommand{\answerTable}[5]{%
\par\vspace{0.15cm}
\noindent
\begin{tabular}{|*{5}{>{\centering\arraybackslash}m{3cm}|}}
\hline
\rule[-0.15cm]{0pt}{0.6cm}\textbf{А} & \rule[-0.15cm]{0pt}{0.6cm}\textbf{Б} & \rule[-0.15cm]{0pt}{0.6cm}\textbf{В} & \rule[-0.15cm]{0pt}{0.6cm}\textbf{Г} & \rule[-0.15cm]{0pt}{0.6cm}\textbf{Д} \\
\hline
\rule[-0.2cm]{0pt}{0.75cm}#1 & \rule[-0.2cm]{0pt}{0.75cm}#2 & \rule[-0.2cm]{0pt}{0.75cm}#3 & \rule[-0.2cm]{0pt}{0.75cm}#4 & \rule[-0.2cm]{0pt}{0.75cm}#5 \\
\hline
\end{tabular}
\par\vspace{0.25cm}
}

\newcommand{\answerTableTall}[5]{%
\par\vspace{0.15cm}
\noindent
\begin{tabular}{|*{5}{>{\centering\arraybackslash}m{3cm}|}}
\hline
\rule[-0.2cm]{0pt}{0.7cm}\textbf{А} & \rule[-0.2cm]{0pt}{0.7cm}\textbf{Б} & \rule[-0.2cm]{0pt}{0.7cm}\textbf{В} & \rule[-0.2cm]{0pt}{0.7cm}\textbf{Г} & \rule[-0.2cm]{0pt}{0.7cm}\textbf{Д} \\
\hline
\rule[-0.45cm]{0pt}{1.2cm}#1 & \rule[-0.45cm]{0pt}{1.2cm}#2 & \rule[-0.45cm]{0pt}{1.2cm}#3 & \rule[-0.45cm]{0pt}{1.2cm}#4 & \rule[-0.45cm]{0pt}{1.2cm}#5 \\
\hline
\end{tabular}
\par\vspace{0.25cm}
}

\newcommand{\instructionBox}[1]{%
\par\vspace{0.3cm}
\begin{tcolorbox}[colback=instrBg, colframe=instrBorder, boxrule=0.6pt, arc=2pt,
    left=8pt, right=8pt, top=4pt, bottom=4pt]
\centering\small\bfseries #1
\end{tcolorbox}
\par\vspace{0.15cm}
}

\newcommand{\sectionTitle}[1]{%
\par\vspace{0.4cm}
\begin{tcolorbox}[colback=titleBg, colframe=titleBg, boxrule=0pt, arc=8pt, halign=center, fontupper=\Large\bfseries\color{mainGreen}]
#1
\end{tcolorbox}\par\vspace{0.3cm}}
\newcommand{\chapterTitle}[1]{\clearpage\sectionTitle{#1}}

\newcommand{\nmtAnswerBox}{%
\par\vspace{0.2cm}\noindent
Відповідь:\ \ 
\begin{tabular}{@{}*{4}{|>{\centering\arraybackslash}p{0.8cm}}|@{\hspace{0.1cm}\textbf{,}\hspace{0.1cm}}*{3}{|>{\centering\arraybackslash}p{0.8cm}}|@{}}
\hline
\rule[-0.2cm]{0pt}{0.7cm} & & & & & & \\
\hline
\end{tabular}
\par\vspace{0.3cm}}

\newcommand{\matchingGrid}{%
    \begingroup
    \renewcommand{\arraystretch}{1.15}
    \setlength{\tabcolsep}{4pt}
    \footnotesize
    \begin{tabular}{r|c|c|c|c|c|}
         \multicolumn{1}{c}{} & \multicolumn{1}{c}{\textbf{А}} & \multicolumn{1}{c}{\textbf{Б}} & \multicolumn{1}{c}{\textbf{В}} & \multicolumn{1}{c}{\textbf{Г}} & \multicolumn{1}{c}{\textbf{Д}} \\ \cline{2-6}
         \textbf{1} & \phantom{X} & \phantom{X} & \phantom{X} & \phantom{X} & \phantom{X} \\ \cline{2-6}
         \textbf{2} & & & & & \\ \cline{2-6}
         \textbf{3} & & & & & \\ \cline{2-6}
    \end{tabular}
    \endgroup
}

\newcounter{zad}
\newcommand{\zadtask}[1]{\stepcounter{zad}\par\vspace{0.3cm}\noindent\textbf{\thezad.}\ \ #1\par\vspace{0.2cm}}
\newcommand{\zadnum}{\stepcounter{zad}\noindent\textbf{\thezad.}\ \ }

\begin{document}
\thispagestyle{empty}
\vspace*{1.2cm}
\begin{center}
{\Large\bfseries\color{mainGreen} РЕАЛЬНИЙ ВАРІАНТ НМТ-2026}\\[0.4cm]
{\Huge\bfseries\color{mainGreen} ВАРІАНТ ЗА 15 ЧЕРВНЯ}\\[0.6cm]
{\large Real Test Notebook з усіма геометричними рисунками}\\[1cm]
{\large\color{mainGreen}\textbf{\href{https://t.me/pvtr2525}{Телеграм: @pvtr2525}}}
\end{center}
\clearpage

\chapterTitle{ТЕСТОВИЙ ЗОШИТ (15 ЧЕРВНЯ)}
\instructionBox{Завдання 1--15 мають по п'ять варіантів відповіді, з яких лише один правильний. Виберіть правильний варіант відповіді й позначте його.}

% ===== №1 =====
\zadtask{Розв'яжіть рівняння $0{,}01x = -1$.}
\answerTable{$-1000$}{$100$}{$-10$}{$-100$}{$-1$}

% ===== №2 =====
\noindent
\begin{minipage}[c]{0.44\textwidth}
\zadnum На діаграмі відображено інформацію про кількість золотих, срібних і бронзових медалей, які вибороли українські спортсмени на олімпійських іграх із 2008 до 2024 року. Якого року українські спортсмени вибороли найбільше медалей?
\end{minipage}%
\hfill
\begin{minipage}[c]{0.54\textwidth}
\centering
\begin{tikzpicture}[x=0.92cm, y=0.28cm, thick]
    % осі та горизонтальна сітка з позначками 0,2,...,14
    \draw[->] (0,0) -- (0,15.2) node[above, font=\scriptsize] {Кількість медалей};
    \draw[->] (0,0) -- (5.6,0) node[below right, font=\scriptsize] {Рік};
    \foreach \y in {0,2,4,6,8,10,12,14} {
        \draw[gray!30, thin] (0,\y) -- (5.3,\y);
        \node[left, font=\scriptsize] at (0,\y) {\y};
    }
    % групи стовпців: золото/срібло/бронза для кожного року
    % 2008: 7,5,15 -> сума 27? масштаб під 14 -> беремо реальні висоти в межах 14
    % значення (золото, срібло, бронза) за реальною діаграмою; найбільша сума - 2008 (22)
    \foreach \i/\year/\g/\s/\b in {0/2008/7/4/11, 1/2012/5/4/10, 2/2016/2/5/4, 3/2020/1/6/12, 4/2024/3/5/4} {
        \pgfmathsetmacro{\bx}{0.55 + \i*1.0}
        \filldraw[fill=gold, draw=black!60, line width=0.3pt] (\bx,0) rectangle ++(0.22,\g);
        \filldraw[fill=silver, draw=black!60, line width=0.3pt] (\bx+0.22,0) rectangle ++(0.22,\s);
        \filldraw[fill=bronze, draw=black!60, line width=0.3pt] (\bx+0.44,0) rectangle ++(0.22,\b);
        \node[below, font=\tiny] at (\bx+0.33,0) {\year};
    }
    % легенда
    \filldraw[fill=gold, draw=black!60] (3.7,13.4) rectangle ++(0.3,0.7);
    \node[right, font=\tiny] at (4.0,13.75) {золото};
    \filldraw[fill=silver, draw=black!60] (3.7,12.2) rectangle ++(0.3,0.7);
    \node[right, font=\tiny] at (4.0,12.55) {срібло};
    \filldraw[fill=bronze, draw=black!60] (3.7,11.0) rectangle ++(0.3,0.7);
    \node[right, font=\tiny] at (4.0,11.35) {бронза};
\end{tikzpicture}
\end{minipage}
\par\vspace{0.2cm}
\answerTable{$2008$}{$2012$}{$2016$}{$2020$}{$2024$}

% ===== №3 =====
\zadtask{Розкладіть на множники вираз $x^2 - 36$.}
\answerTable{$(x-18)(x+18)$}{$(x-6)^2$}{$(x-6)(x+6)$}{$(x+6)^2$}{$(x-18)^2$}

% ===== №4 =====
\noindent
\begin{minipage}[c]{0.60\textwidth}
\zadnum Розгляньте схематичне зображення стелажа для книг і декору. Зображення складається з восьми рівних квадратів. Визначте площу одного такого квадрата, якщо висота й ширина стелажа дорівнюють $160$~см і $80$~см відповідно.
\end{minipage}%
\hfill
\begin{minipage}[c]{0.38\textwidth}
\centering
\begin{tikzpicture}[scale=0.7, thick, line join=round]
    % стелаж: 2 стовпці x 4 ряди квадратів (сторона 40 -> W=80, H=160)
    \draw (0,0) rectangle (2,8);
    \draw (1,0) -- (1,8);
    \foreach \y in {2,4,6} {\draw (0,\y) -- (2,\y);}
    \node[below, font=\scriptsize] at (1,0) {$80$~см};
    \node[left, font=\scriptsize, rotate=90, anchor=south] at (0,4) {$160$~см};
\end{tikzpicture}
\end{minipage}
\par\vspace{0.2cm}
\answerTable{$800$~см$^2$}{$1200$~см$^2$}{$2000$~см$^2$}{$160$~см$^2$}{$1600$~см$^2$}

% ===== №5 =====
\zadtask{Позначте назву многогранника, який має дві основи й дев'ять ребер.
\vspace{0.1cm}
\begin{enumerate}[label=\textbf{\Alph*}, leftmargin=2.8em, itemsep=0.1cm]
    \item чотирикутна призма
    \item трикутна призма
    \item дев'ятикутна призма
    \item чотирикутна піраміда
    \item трикутна піраміда
\end{enumerate}
\vspace{0.1cm}}

% ===== №6 =====
\zadtask{Обчисліть $\log_4 2 + \log_4 32$.}
\answerTable{$\log_4 34$}{$3$}{$2$}{$4$}{$16$}

% ===== №7 =====
\zadtask{Розв'яжіть нерівність $5^{x-1} \leqslant \dfrac{1}{25}$.}
\answerTable{$(-\infty; -1]$}{$(-\infty; -4]$}{$(-\infty; -3]$}{$[1; +\infty)$}{$[3; +\infty)$}

% ===== №8 =====
\noindent
\begin{minipage}[c]{0.58\textwidth}
\zadnum На рисунку зображено графік функції $y = f(x)$, визначеної на проміжку $[-3;\, 3]$. На якому з наведених проміжків ця функція спадає?
\end{minipage}%
\hfill
\begin{minipage}[c]{0.40\textwidth}
\centering
\begin{tikzpicture}[scale=0.5, thick]
    \draw[step=1cm, gray!40, thin] (-3,-2) grid (4,5);
    \draw[->] (-3.3,0) -- (4.3,0) node[right, font=\scriptsize] {$x$};
    \draw[->] (0,-2.3) -- (0,5.3) node[above, font=\scriptsize] {$y$};
    \node[below left, font=\scriptsize] at (0,0) {$0$};
    \node[below, font=\scriptsize] at (-3,0) {$-3$};
    \node[below, font=\scriptsize] at (1,0) {$1$};
    \node[below, font=\scriptsize] at (3,0) {$3$};
    \node[left, font=\scriptsize] at (0,1) {$1$};
    \draw[mainGreen, very thick, smooth] plot coordinates {(-3,-2) (-1,0.5)(0,3.2) (1,4) (2,3.2) (3,1)};
    \node[right, font=\scriptsize] at (1.5,4.6) {$y=f(x)$};
\end{tikzpicture}
\end{minipage}
\par\vspace{0.2cm}
\answerTable{$[-3; 1]$}{$[-3; 4]$}{$[2; 4]$}{$[1; 3]$}{$[-3; 3]$}

% ===== №9 =====
\zadtask{Які твердження щодо кутів на площині є правильними?
\vspace{0.1cm}
\begin{enumerate}[label=\textbf{\Roman*.}, align=left, leftmargin=2.6em, labelsep=0.5em, itemsep=0.08cm, topsep=0.06cm]
\item Якщо два прямі кути мають спільну сторону, то інші сторони цих кутів лежать на одній прямій.
\item Якщо два кути мають спільну вершину, то вони вертикальні.
\item Якщо два кути мають спільну сторону, то вони суміжні.
\end{enumerate}
\vspace{0.1cm}}
\answerTable{лише I}{лише II}{лише III}{лише I та II}{лише II та III}

% ===== №10 =====
\zadtask{Знайдіть первісну для функції $f(x) = \dfrac{1}{3x}$.}
\answerTableTall{$F(x) = \ln|3x|$}{$F(x) = \dfrac{1}{3}\ln|x|$}{$F(x) = 3\ln|x|$}{$F(x) = -\dfrac{1}{3x^2}$}{$F(x) = \dfrac{2}{3x^2}$}
\newpage

% ===== №11 =====
\noindent
\begin{minipage}[c]{0.54\textwidth}
\zadnum На круговій діаграмі відображено розподіл музичних треків за жанрами на платівці. Програвач відтворює один випадковий трек із цієї платівки. Знайдіть ймовірність того, що першим прозвучить трек у жанрі рок.
\end{minipage}%
\hfill
\begin{minipage}[c]{0.44\textwidth}
\centering
\begin{tikzpicture}[scale=1.0]
    \def\R{1.9}
    % значення: рок 20, поп 25, джаз 11, диско 14; сума 70
    % кути: рок 20/70*360=102.857; поп 25/70*360=128.571; джаз 11/70=56.571; диско 14/70=72
    % старт зверху (90) за годинниковою
    \definecolor{crok}{RGB}{80,180,220}
    \definecolor{cpop}{RGB}{235,90,95}
    \definecolor{cjazz}{RGB}{95,190,120}
    \definecolor{cdisco}{RGB}{245,200,70}
    % рок: 90 -> 90-102.857=-12.857
    \filldraw[fill=crok, draw=white, line width=1pt] (0,0) -- (90:\R) arc (90:-12.857:\R) -- cycle;
    \node at (38:1.1) {$20$};
    % поп: -12.857 -> -141.428
    \filldraw[fill=cpop, draw=white, line width=1pt] (0,0) -- (-12.857:\R) arc (-12.857:-141.428:\R) -- cycle;
    \node at (-77:1.1) {$25$};
    % джаз: -141.428 -> -198 (=162)
    \filldraw[fill=cjazz, draw=white, line width=1pt] (0,0) -- (-141.428:\R) arc (-141.428:-198:\R) -- cycle;
    \node at (-170:1.15) {$11$};
    % диско: -198 -> -270 (=90)
    \filldraw[fill=cdisco, draw=white, line width=1pt] (0,0) -- (-198:\R) arc (-198:-270:\R) -- cycle;
    \node at (-234:1.15) {$14$};
    % легенда
    \begin{scope}[shift={(2.5,1.3)}]
        \filldraw[fill=crok, draw=black!40] (0,0) rectangle ++(0.3,0.3); \node[right, font=\scriptsize] at (0.35,0.15){рок};
        \filldraw[fill=cpop, draw=black!40] (0,-0.55) rectangle ++(0.3,0.3); \node[right, font=\scriptsize] at (0.35,-0.4){поп};
        \filldraw[fill=cjazz, draw=black!40] (0,-1.1) rectangle ++(0.3,0.3); \node[right, font=\scriptsize] at (0.35,-0.95){джаз};
        \filldraw[fill=cdisco, draw=black!40] (0,-1.65) rectangle ++(0.3,0.3); \node[right, font=\scriptsize] at (0.35,-1.5){диско};
    \end{scope}
\end{tikzpicture}
\end{minipage}
\par\vspace{0.2cm}
\answerTableTall{$\dfrac{2}{5}$}{$\dfrac{5}{14}$}{$\dfrac{1}{20}$}{$\dfrac{1}{4}$}{$\dfrac{2}{7}$}

% ===== №12 =====
\zadtask{Визначте координати вектора $\vec{c} = \vec{a} - \vec{b}$, якщо $\vec{a}(3;\, 5;\, 7)$ і $\vec{b}(2;\, -4;\, 8)$.}
\answerTable{$(5; 1; 15)$}{$(-4; -9; 1)$}{$(-1; -9; 1)$}{$(1; 1; -1)$}{$(1; 9; -1)$}

% ===== №13 =====
\zadtask{$\sin^2 120^\circ \cdot \cos 120^\circ = $}
\answerTableTall{$-\dfrac{1}{8}$}{$-\dfrac{\sqrt{3}}{8}$}{$-\dfrac{3}{8}$}{$-\dfrac{3\sqrt{3}}{8}$}{$-\dfrac{\sqrt{2}}{4}$}

% ===== №14 =====
\noindent
\begin{minipage}[c]{0.60\textwidth}
\zadnum У рівнобедреному трикутнику $ABC$ з точки $K$, що лежить на бічній стороні $BC$, до бічної сторони $AB$ проведено перпендикуляр $KM$ (див. рисунок). Визначте площу трикутника $ABC$, якщо $KM = 6\sqrt{3}$~см, $AK : MB = 2 : 3$, $\angle B = 30^\circ$.
\end{minipage}%
\hfill
\begin{minipage}[c]{0.38\textwidth}
\centering
\begin{tikzpicture}[scale=0.72, thick, line join=round]
    \coordinate (A) at (0,0);
    \coordinate (C) at (3.2,0);
    \coordinate (B) at (1.05,4.2);
    % M на AB, K на BC
    \coordinate (M) at ($(A)!0.55!(B)$);
    \coordinate (K) at ($(B)!0.55!(C)$);
    \draw (A) -- (B) -- (C) -- cycle;
    \draw (K) -- (M);
    \draw (A) -- (K);
    \pic[draw, angle radius=0.18cm] {right angle = K--M--A};
    \pic[draw, angle radius=0.5cm, pic text={$30^\circ$}, angle eccentricity=1.55, font=\scriptsize] {angle = A--B--C};
    \node[below left] at (A) {$A$};
    \node[above] at (B) {$B$};
    \node[below right] at (C) {$C$};
    \node[left] at (M) {$M$};
    \node[right] at (K) {$K$};
\end{tikzpicture}
\end{minipage}
\par\vspace{0.2cm}
\answerTable{$288\sqrt{3}$~см$^2$}{$144\sqrt{3}$~см$^2$}{$72\sqrt{3}$~см$^2$}{$288$~см$^2$}{$144$~см$^2$}

% ===== №15 =====
\zadtask{Якому проміжку належить \textit{більший} корінь рівняння $\sqrt{\dfrac{2x^2 + 7x + 10}{2}} = 4$?}
\answerTable{$(-\infty; -4)$}{$[-4; -1)$}{$[-1; 1)$}{$[1; 4)$}{$[4; +\infty)$}
\newpage
\instructionBox{У завданнях 16--18 до кожного з трьох рядків інформації, позначених цифрами, доберіть один правильний, на Вашу думку, варіант, позначений буквою.}

% ===== №16 =====
\zadtask{Узгодьте вираз (1--3) з тотожно рівним йому виразом (А--Д), якщо $a > 0$, $a \neq 1$.\par
\vspace{0.3cm}
\noindent
\begin{minipage}[t]{0.42\textwidth}
\textit{Вираз}\par\vspace{0.15cm}
\textbf{1} \quad $\sqrt{a^2}$\\[0.35cm]
\textbf{2} \quad $a^3 : a^4$\\[0.35cm]
\textbf{3} \quad $\dfrac{a^2 - a}{1 - a}$
\end{minipage}%
\hfill
\begin{minipage}[t]{0.3\textwidth}
\textit{Тотожно рівний вираз}\par\vspace{0.15cm}
\textbf{А} \quad $a$\\[0.15cm]
\textbf{Б} \quad $2a$\\[0.15cm]
\textbf{В} \quad $\dfrac{1}{a}$\\[0.15cm]
\textbf{Г} \quad $-a$\\[0.15cm]
\textbf{Д} \quad $|a|$
\end{minipage}%
\hfill
\begin{minipage}[t]{0.2\textwidth}
\vspace{0.4cm}
\begin{flushright}\matchingGrid\end{flushright}
\end{minipage}}

% ===== №17 =====
\zadtask{До кожного початку речення (1--3) доберіть його закінчення (А--Д) так, щоб утворилося правильне твердження.\par
\vspace{0.3cm}
\noindent
\begin{minipage}[t]{0.50\textwidth}
\textit{Початок речення}\par\vspace{0.15cm}
\textbf{1} \quad Графіки функцій $y = 2^x$ та $y = 0{,}5^x$ перетинаються в точці\\[0.25cm]
\textbf{2} \quad Пряма $y = x - 1$ перетинає вісь $x$ у точці\\[0.25cm]
\textbf{3} \quad Графік функції $y = (x + 2)^2$ дотикається осі $x$ у точці
\end{minipage}%
\hfill
\begin{minipage}[t]{0.28\textwidth}
\textit{Закінчення речення}\par\vspace{0.15cm}
\textbf{А} \quad $(1; 0)$.\\[0.15cm]
\textbf{Б} \quad $(0; 1)$.\\[0.15cm]
\textbf{В} \quad $(2; 0)$.\\[0.15cm]
\textbf{Г} \quad $(-2; 0)$.\\[0.15cm]
\textbf{Д} \quad $(-1; 0)$.
\end{minipage}%
\hfill
\begin{minipage}[t]{0.18\textwidth}
\vspace{0.4cm}
\begin{flushright}\matchingGrid\end{flushright}
\end{minipage}}

% ===== №18 =====
\zadtask{На рисунку точки $K$, $L$, $M$ та $N$~--- середини сторін $AB$, $BC$, $CD$ та $AD$ прямокутника $ABCD$ відповідно. $AB = 12$~см, $BC = 16$~см. Точка $P$~--- середина відрізка $KN$. Установіть відповідність між відрізком (1--3) та його довжиною (А--Д).
\vspace{0.3cm}
\begin{center}
\begin{tikzpicture}[scale=0.55, thick, line join=round]
    \coordinate (A) at (0,0);
    \coordinate (B) at (0,12);
    \coordinate (C) at (16,12);
    \coordinate (D) at (16,0);
    \coordinate (K) at (0,6);
    \coordinate (L) at (8,12);
    \coordinate (M) at (16,6);
    \coordinate (N) at (8,0);
    \coordinate (P) at (4,3);
    \draw (A) -- (B) -- (C) -- (D) -- cycle;
    \draw (K) -- (L) -- (M) -- (N) -- cycle;
    \fill (P) circle (3pt);
    \draw (K)--(N);
    \path (K) -- (N) node[midway, sloped] {\footnotesize $|$};
    \node[below left] at (A) {$A$};
    \node[above left] at (B) {$B$};
    \node[above right] at (C) {$C$};
    \node[below right] at (D) {$D$};
    \node[left] at (K) {$K$};
    \node[above] at (L) {$L$};
    \node[right] at (M) {$M$};
    \node[below] at (N) {$N$};
    \node[above left, font=\scriptsize] at (P) {$P$};
\end{tikzpicture}
\end{center}
\vspace{0.2cm}
\noindent
\begin{minipage}[t]{0.52\textwidth}
\textit{Відрізок}\par\vspace{0.15cm}
\textbf{1} \quad відстань від точки $P$ до сторони $BC$\\[0.25cm]
\textbf{2} \quad відстань від точки $P$ до середини $MN$\\[0.25cm]
\textbf{3} \quad відстань від точки $P$ до вершини $C$
\end{minipage}
\hfill
\begin{minipage}[t]{0.26\textwidth}
\textit{Довжина відрізка}\par\vspace{0.15cm}
\textbf{А} \quad $8$~см\\[0.15cm]
\textbf{Б} \quad $9$~см\\[0.15cm]
\textbf{В} \quad $12$~см\\[0.15cm]
\textbf{Г} \quad $15$~см\\[0.15cm]
\textbf{Д} \quad $16$~см\\[0.15cm]
\end{minipage}
\hfill
\begin{minipage}[t]{0.18\textwidth}
\vspace{0.4cm}
\begin{flushright}\matchingGrid\end{flushright}
\end{minipage}}
\newpage
\instructionBox{Розв'яжіть завдання 19--22. Одержані числові відповіді запишіть у спеціально відведеному місці. Відповідь записуйте лише десятковим дробом, урахувавши положення коми. Знак «мінус» записуйте перед першою цифрою числа.}

% ===== №19 =====
\zadtask{Визначте найбільше значення функції $y = 8\sqrt{x} - x + 4$, якщо $x \in [0;\, 36]$.}
\nmtAnswerBox

% ===== №20 =====
\zadtask{У грушевому саду ростуть $200$ груш п'яти різних сортів. У таблиці наведено інформацію про назви сортів груш і частку (у~\%) дерев цих сортів від загальної кількості груш у саду.
\vspace{0.2cm}

\noindent\begin{center}
\renewcommand{\arraystretch}{1.3}
\begin{tabular}{|>{\centering\arraybackslash}m{2.2cm}|c|c|>{\centering\arraybackslash}m{2cm}|>{\centering\arraybackslash}m{2cm}|c|}
\hline
\textbf{Сорт груш} & Дюшес & Вільямс & Улюблениця Клаппа & Лісова Красуня & Бланка \\
\hline
\textbf{Частка (\%) дерев} & $21$ & $11$ & $38$ & $n$ & $10$ \\
\hline
\end{tabular}
\end{center}
\vspace{0.2cm}

У таблиці $n$~--- невідоме число. За всі груші сорту Лісова Красуня, реалізовані на оптовому ринку торік, виручили $52\,000$~грн. Цього року запланували за всі груші сорту Лісова Красуня вторгувати на $10\,\%$ більше гривень, ніж торік. Скільки гривень у середньому планують виручити за груші з одного дерева цього сорту?}
\nmtAnswerBox

% ===== №21 =====
\noindent
\begin{minipage}[c]{0.56\textwidth}
\zadnum На рисунку зображено циліндр і конус, які мають рівні основи й об'єми. Визначте площу $S$ (у см$^2$) бічної поверхні циліндра, якщо об'єм конуса дорівнює $810\pi$~см$^3$, а діаметр його основи~--- $18$~см. У відповідь запишіть значення $\dfrac{S}{\pi}$.
\end{minipage}%
\hfill
\begin{minipage}[c]{0.42\textwidth}
\centering
\begin{tikzpicture}[scale=0.6, thick, line join=round]
    % циліндр
    \draw (0,3) ellipse (1.2 and 0.35);
    \draw[dashed] (-1.2,0) arc (180:0:1.2 and 0.35);
    \draw (-1.2,0) arc (180:360:1.2 and 0.35);
    \draw (-1.2,0) -- (-1.2,3);
    \draw (1.2,0) -- (1.2,3);
    \fill (0,0) circle (1.2pt);
    % конус
    \begin{scope}[shift={(4.2,0)}]
        \draw[dashed] (-1.2,0) arc (180:0:1.2 and 0.35);
        \draw (-1.2,0) arc (180:360:1.2 and 0.35);
        \draw (-1.2,0) -- (0,4.2);
        \draw (1.2,0) -- (0,4.2);
        \fill (0,0) circle (1.2pt);
    \end{scope}
\end{tikzpicture}
\end{minipage}
\par\vspace{0.2cm}
\nmtAnswerBox

% ===== №22 =====
\zadtask{Визначте \textit{суму} всіх цілих значень $a$, за кожного з яких добуток різних коренів рівняння
\[
\log_3\!\big(|x - 2| + 2a - 11\big) - 2 = 0
\]
є більшим за $-40$.}
\nmtAnswerBox

\end{document}